% Auto-generated by generate_tables.py — DO NOT EDIT

% Table 1: γ spectral overlap across the Pythia suite
\begin{table}[t]
  \centering
  \caption{LoRA adapter stable rank ($srank$) and subspace-overlap bonus factors across the Pythia model suite. \textbf{bonus} = measured overlap divided by the random-subspace baseline $k/d$; values $> 1$ indicate alignment with the base weight's top-$k$ singular subspace. $k=5$ is the raw probe; $k=srank$ is the natural-rank probe. DPO = Direct Preference Optimization on \texttt{Anthropic/hh-rlhf}; CLM = next-token-prediction control on the same data. All runs: LoRA rank 128, $\alpha = 256$, 800 steps, fp16.}
  \label{tab:gamma-scaling}
  \small
  \begin{tabular}{lrrrrrr}
    \toprule
    Model & $d_{\text{model}}$ & layers & $srank$ & $\text{bonus}_R(5)$ & $\text{bonus}_R(srank)$ & $\text{bonus}_L(5)$ \\
    \midrule
    70M DPO & 512 & 6 & 3.93 & 5.12 & 5.50 & 1.29 \\
    160M DPO & 768 & 12 & 3.51 & 5.41 & 9.20 & 0.75 \\
    410M DPO & 1024 & 24 & 3.92 & 5.06 & 7.80 & 1.47 \\
    410M CLM & 1024 & 24 & 3.45 & 5.22 & 9.21 & 1.42 \\
    1B DPO (seed 42) & 2048 & 16 & 3.13 & 3.24 & 5.38 & 3.27 \\
    1B DPO (seed 117) & 2048 & 16 & 3.36 & 3.49 & 5.37 & 2.83 \\
    1B CLM (seed 42) & 2048 & 16 & 2.89 & 4.15 & 8.23 & 3.27 \\
    1B CLM (seed 117) & 2048 & 16 & 2.83 & 4.02 & 7.36 & 3.15 \\
    \bottomrule
  \end{tabular}
\end{table}

% Table 2: BitFit-DPO gauge theory falsification test
\begin{table}[t]
  \centering
  \caption{Bias-only fine-tuning (BitFit) of Pythia-410M under DPO. With $271,360$ trainable bias parameters ($\approx 0.07\%$ of total), the adapter saturates at a loss far above the LoRA-DPO endpoint ($\approx 0.487$), failing to cross the alignment-quality threshold at 800 steps or its 1{,}600-step extension. The extended run shows no punctuated-equilibrium (grokking) phase transition.}
  \label{tab:bitfit}
  \small
  \begin{tabular}{lrr}
    \toprule
    Quantity & BitFit (800 steps) & BitFit ext. (1560 steps) \\
    \midrule
    Initial loss            & 0.693 & 0.693 \\
    Final loss (last-10 mean) & 0.660 & 0.607 \\
    Min loss observed       & 0.595 & 0.567 \\
    Drop from initial (nats)& 0.033 & 0.086 \\
    Verdict                 & \texttt{gauge\_dead} & \texttt{gauge\_confirmed\_dead} \\
    \bottomrule
  \end{tabular}
\end{table}

% Table 3: Scaling-form fit comparison
\begin{table}[t]
  \centering
  \caption{Least-squares fit of four candidate scaling forms to the measured DPO stable ranks at four model scales. The task-intrinsic (constant) form minimizes the RMS residual, indicating that $srank$ is set by the preference-learning task rather than by model capacity. All values in stable-rank units.}
  \label{tab:scaling-fit}
  \small
  \begin{tabular}{lccccc}
    \toprule
    \multirow{2}{*}{Model} & \multirow{2}{*}{$d$} & \multirow{2}{*}{measured} & \multicolumn{3}{c}{predicted by best fit} \\
    \cmidrule(lr){4-6}
     & & & task-intrinsic & $c/\sqrt{d}$ & $c/\sqrt[3]{d}$ \\
    \midrule
    70M & 512 & 3.93 & 3.62 & 4.61 & 4.33 \\
    160M & 768 & 3.51 & 3.62 & 3.76 & 3.78 \\
    410M & 1024 & 3.92 & 3.62 & 3.26 & 3.44 \\
    1B & 2048 & 3.13 & 3.62 & 2.30 & 2.73 \\
    \midrule
    RMS residual & & & \textbf{0.330} & 0.644 & 0.399 \\
    Best-fit constant & & & 3.62 & 104.3 & 34.7 \\
    \bottomrule
  \end{tabular}
\end{table}

% Table 4: Bias-theory (diagonal-gain) autopsy
\begin{table}[t]
  \centering
  \caption{Decomposition of LoRA adapters on \texttt{attention.query\_key\_value} into diagonal-gain components. We fit $\Delta W \approx W \mathrm{diag}(g)$ (right-multiplicative, equivalent to LayerNorm-$\gamma$ modulation) and $\Delta W \approx \mathrm{diag}(g) W$ (left-multiplicative, output-channel gain), and report the Frobenius-squared residual fraction. 99.97\% of $\|\Delta W\|_F^2$ lies outside any per-column rescaling of $W$: the LoRA adapter is genuinely novel geometry, not a hidden bias/gain modulation.}
  \label{tab:autopsy}
  \small
  \begin{tabular}{lcc}
    \toprule
    Run & Residual fraction (right-mult) & Residual fraction (left-mult) \\
    \midrule
    v1 DPO $r{=}16$  & 99.97\% & 99.89\% \\
    v2 DPO $r{=}128$ & 99.97\% & 99.83\% \\
    v3 CLM $r{=}128$ & 99.97\% & 99.82\% \\
    \bottomrule
  \end{tabular}
\end{table}

% Table 5: Behavior-geometry correlation (T1.2)
\begin{table}[t]
  \centering
  \caption{Behavior-geometry correlation across Pythia DPO checkpoints. For each checkpoint, we compute reward margin $= \beta\cdot[\log\pi_\theta(y_{\text{win}}|x)/\pi_{\text{ref}}(y_{\text{win}}|x) - \log\pi_\theta(y_{\text{los}}|x)/\pi_{\text{ref}}(y_{\text{los}}|x)]$ and sequence-level KL-to-base $= n^{-1}\sum_t[\log\pi_\theta(y_t|y_{<t}) - \log\pi_{\text{ref}}(y_t|y_{<t})]$ on the \texttt{Anthropic/hh-rlhf} test split ($n=500$, $\beta=0.1$). Geometry columns (srank, $\gamma$) are from prior spectral analysis. With $n=5$ checkpoints, correlations are suggestive, not definitive.}
  \label{tab:behavior-geometry}
  \small
  \begin{tabular}{lrrrrrr}
    \toprule
    Checkpoint & $srank$ & $\gamma_{R}(k{=}5)$ & rew.\ margin & $\pm$ SE & KL-to-base & $\pm$ SE \\
    \midrule
    70m (s42) & 3.93 & 5.12 & 0.0517 & 0.0237 & -0.0553 & 0.0061 \\
    160m (s42) & 3.51 & 5.41 & 0.0287 & 0.0218 & -0.0637 & 0.0059 \\
    410m (s42) & 3.92 & 5.06 & 0.0739 & 0.0319 & -0.1206 & 0.0085 \\
    1b (s42) & 3.13 & 3.24 & 0.0561 & 0.0259 & -0.0975 & 0.0071 \\
    1b (s117) & 3.36 & 3.49 & 0.0618 & 0.0265 & -0.1083 & 0.0078 \\
    \midrule
    \multicolumn{2}{l}{srank vs rew.~margin} & \multicolumn{2}{l}{Pearson: 0.22 [-0.15,1.00]} & \multicolumn{2}{l}{Spearman: -0.10} \\
    \multicolumn{2}{l}{$\gamma$ vs rew.~margin} & \multicolumn{2}{l}{Pearson: -0.35 [-1.00,-0.79]} & \multicolumn{2}{l}{Spearman: -0.60} \\
    \multicolumn{2}{l}{srank vs KL} & \multicolumn{2}{l}{Pearson: 0.18 [-0.24,0.96]} & \multicolumn{2}{l}{Spearman: 0.30} \\
    \multicolumn{2}{l}{$\gamma$ vs KL} & \multicolumn{2}{l}{Pearson: 0.49 [0.29,0.96]} & \multicolumn{2}{l}{Spearman: 0.50} \\
    \bottomrule
  \end{tabular}
\end{table}
